\documentclass[12pt]{article}
\usepackage[margin=1in]{geometry}
\usepackage{enumitem}
\usepackage{titlesec}
\usepackage{parskip}
\usepackage{hyperref}
\usepackage{fontenc}

\title{\textbf{Contemporary Comparative Politics \& Political Institutions}\\
\large A Supplementary Reading List of Core Books Published After 2000\\
(with a political-economy rooting, and a note on UC Berkeley affiliations)}
\author{}
\date{}

\begin{document}

\maketitle
\tableofcontents
\newpage

%--------------------------------------------------
\section{Scope and Purpose}
%--------------------------------------------------

\begin{itemize}[leftmargin=*, itemsep=4pt]
    \item This is a curated (not exhaustive) supplement to the core comp-exam list, gathering major books in comparative politics and political institutions---including work with an explicit political-economy rooting---published after 2000, weighted toward the more recent literature.
    \item Entries are grouped by the field's main clusters. Each carries a one-line note on why it matters or how it fits the contemporary debate.
    \item The final section flags authors who were on the UC Berkeley faculty during 2008--2012, per the specific interest in identifying that lineage. A caution about mis-attribution (Dunning) is included there.
\end{itemize}

%--------------------------------------------------
\section{Institutions \& Comparative Political Economy}
%--------------------------------------------------

\begin{itemize}[leftmargin=*, itemsep=4pt]
    \item \textbf{Peter Hall \& David Soskice (eds.), \textit{Varieties of Capitalism} (2001)} --- the foundational post-2000 framework for institutional variation among rich democracies; most of CPE since responds to it.
    \item \textbf{Kathleen Thelen, \textit{How Institutions Evolve} (2004)} and \textbf{\textit{Varieties of Liberalization and the New Politics of Social Solidarity} (2014)} --- the reference works on gradual institutional change; the 2014 book is the more contemporary statement.
    \item \textbf{Paul Pierson, \textit{Politics in Time: History, Institutions, and Social Analysis} (2004)} --- the canonical treatment of path dependence and temporality in institutional analysis. (Pierson $\rightarrow$ Berkeley; see \S8.)
    \item \textbf{Torben Iversen \& David Soskice, \textit{Democracy and Prosperity} (2019)} --- a major recent synthesis arguing advanced capitalism and democracy are mutually reinforcing; a useful ``state of the field'' capstone.
    \item \textbf{Daron Acemoglu \& James Robinson, \textit{Economic Origins of Dictatorship and Democracy} (2006), \textit{Why Nations Fail} (2012), \textit{The Narrow Corridor} (2019)} --- the dominant institutionalist--political-economy sequence on regimes and development.
    \item \textbf{Douglass North, John Wallis \& Barry Weingast, \textit{Violence and Social Orders} (2009)} --- ``limited vs.\ open access orders''; an ambitious institutional theory of the state and violence.
    \item \textbf{Timothy Besley \& Torsten Persson, \textit{Pillars of Prosperity} (2011)} --- state capacity (fiscal and legal) as the core institutional variable.
\end{itemize}

%--------------------------------------------------
\section{Regimes: Democracy, Authoritarianism, Backsliding}
%--------------------------------------------------

\begin{itemize}[leftmargin=*, itemsep=4pt]
    \item \textbf{Steven Levitsky \& Lucan Way, \textit{Competitive Authoritarianism} (2010)} --- defined the hybrid-regime research agenda; still the central reference. Their \textit{Revolution and Dictatorship} (2022) extends it.
    \item \textbf{Milan Svolik, \textit{The Politics of Authoritarian Rule} (2012)} --- the leading formal-institutional account of how dictatorships manage elite and mass threats.
    \item \textbf{Jennifer Gandhi, \textit{Political Institutions under Dictatorship} (2008)} and \textbf{Beatriz Magaloni, \textit{Voting for Autocracy} (2006)} --- legislatures, parties, and elections \emph{inside} authoritarian regimes.
    \item \textbf{Barbara Geddes, Joseph Wright \& Erica Frantz, \textit{How Dictatorships Work} (2018)} --- the synthesis of the large-$N$ authoritarian-regimes program.
    \item \textbf{Carles Boix, \textit{Democracy and Redistribution} (2003)} and \textbf{\textit{Political Order and Inequality} (2015)} --- inequality-based theories of regime origins and order.
    \item \textbf{Daniel Ziblatt, \textit{Conservative Parties and the Birth of Democracy} (2017)} --- why democratization succeeds or fails depending on whether the old elite can organize; excellent institutional history.
    \item \textbf{Stephan Haggard \& Robert Kaufman, \textit{Development, Democracy, and Welfare States} (2008)} and \textbf{\textit{Backsliding} (2021)} --- the latter is among the best recent treatments of democratic erosion.
    \item \textbf{Dan Slater, \textit{Ordering Power} (2010)} --- authoritarian durability and state-building in Southeast Asia; a modern comparative-historical classic.
\end{itemize}

%--------------------------------------------------
\section{Parties, Elections, and Distributive Politics}
%--------------------------------------------------

\begin{itemize}[leftmargin=*, itemsep=4pt]
    \item \textbf{Herbert Kitschelt \& Steven Wilkinson (eds.), \textit{Patrons, Clients, and Policies} (2007)} and \textbf{Susan Stokes, Thad Dunning, Marcelo Nazareno \& Valeria Brusco, \textit{Brokers, Voters, and Clientelism} (2013)} --- the anchors of the clientelism / distributive-politics literature.
    \item \textbf{Kanchan Chandra, \textit{Why Ethnic Parties Succeed} (2004)} --- patronage and ethnic mobilization.
    \item \textbf{Pradeep Chhibber \& Rahul Verma, \textit{Ideology and Identity: The Changing Party Systems of India} (2018)} --- recent, and a Berkeley author (see \S8).
\end{itemize}

%--------------------------------------------------
\section{Violence, Order, and State Formation}
%--------------------------------------------------

\begin{itemize}[leftmargin=*, itemsep=4pt]
    \item \textbf{Stathis Kalyvas, \textit{The Logic of Violence in Civil War} (2006)} --- the defining book on wartime violence; reoriented an entire subfield.
    \item \textbf{James C. Scott, \textit{The Art of Not Being Governed} (2009)} --- statelessness and the limits of state reach (pairs with his 1998 \textit{Seeing Like a State}).
\end{itemize}

%--------------------------------------------------
\section{Comparative Political Economy of the Advanced / Asian Economies}
%--------------------------------------------------

\begin{itemize}[leftmargin=*, itemsep=4pt]
    \item \textbf{Steven K. Vogel, \textit{Japan Remodeled} (2006)} and \textbf{\textit{Marketcraft: How Governments Make Markets Work} (2018)} --- markets as governed institutions. (Vogel is at Berkeley; see \S8.)
    \item \textbf{Jonah Levy (ed.), \textit{The State after Statism} (2006)} --- reconfiguration of the state's economic role across rich democracies. (Levy is at Berkeley; see \S8.)
\end{itemize}

%--------------------------------------------------
\section{The UC Berkeley 2008--2012 Connection}
%--------------------------------------------------

\begin{itemize}[leftmargin=*, itemsep=4pt]
    \item \textbf{Paul Pierson} --- moved to Berkeley in 2007 and was on the faculty throughout 2008--2012. Beyond \textit{Politics in Time}, his co-authored \textit{Winner-Take-All Politics} (2010, with Jacob Hacker) is a signature contemporary political-economy book on institutions and inequality.
    \item \textbf{Steven K. Vogel} --- Berkeley faculty across that period; \textit{Japan Remodeled} (2006) and \textit{Marketcraft} (2018).
    \item \textbf{Pradeep Chhibber} --- Berkeley faculty; parties and Indian politics (\textit{Ideology and Identity}, 2018).
    \item \textbf{Jonah Levy}, \textbf{T.\,J.\ Pempel} (\textit{Regime Shift}), and \textbf{Ruth Berins Collier} were also comparative / political-economy faculty there in those years, if the net is widened.
    \item \textbf{Caution against mis-attribution:} \textbf{Thad Dunning} (co-author of \textit{Brokers, Voters, and Clientelism}) is \emph{now} at Berkeley but did not join until $\sim$2014---during 2008--2012 he was at Yale. If the ``Berkeley-in-that-window'' criterion is strict, lean on Pierson, Vogel, Chhibber, and Levy.
\end{itemize}

%--------------------------------------------------
\section{Optional Extensions}
%--------------------------------------------------

\begin{itemize}[leftmargin=*, itemsep=4pt]
    \item \textbf{If a tighter core is wanted:} Hall \& Soskice (2001), Thelen (2004), Pierson (2004), Acemoglu \& Robinson (2006/2012), North--Wallis--Weingast (2009), Levitsky \& Way (2010), Svolik (2012), Boix (2003), Ziblatt (2017), Slater (2010), Kalyvas (2006), Stokes et al.\ (2013) make a defensible twelve.
    \item \textbf{If comparative methods should be added:} John Gerring, \textit{Case Study Research} (2007); Thad Dunning, \textit{Natural Experiments in the Social Sciences} (2012); and David Collier's work on concepts and comparison.
\end{itemize}

\end{document}
